\begin{abstract}
% v21, 2026-09-08. 205 words / 10 sentences. FINAL for this pass.
% Flesch 50.8 (above all 147 measured ACL-venue abstracts, median 1.2, max 40.7).
% Density 2.93 numerals per 100 words = 91st percentile (corpus median 0.49, p75 1.36).
% 3 result percentages, all traced to results_FINAL.md and all reported in results_v2.tex.
% Register: 0 em dashes, colons, semicolons, spelled fractions, imperatives, "we show",
% significance adverbs, not-X-but-Y, sentence-final prepositions.
% Only benchmark still missed: length, 205 against a venue median of 161 (max observed 232).
% Every shortening attempt below ~200 cost a beat we decided to keep.
Do language models know when a good answer should be left alone? They are increasingly
handed work their users cannot judge, and a user who cannot name the fault falls back on
undirected requests that carry no information about what to change. Self-correction research
asks whether a model can repair a wrong response, and finds that it can with a specific
critique and largely cannot without one. That leaves unexamined the case where the work is
already sufficient, as 88\% of first drafts here are, and gets revised anyway because the
user cannot tell. In this work, we run five-turn revision conversations with 6 models on 40
tasks across 5 domains, contrasting undirected requests with a single targeted critique.
Most replies contain no revision at all, with models restating the work or declining while
presenting it as compliance. Revisions that change quality make it worse 70\% of the time.
Applied to work that was already sufficient, 27\% of them leave it insufficient. For every
model the first draft rates highest, and blind readers show no reliable preference between
it and the last. These results indicate that undirected revision lowers the quality of work
that was already sufficient, and that naming the fault reverses the effect.
\end{abstract}
